% Generated draft: requires mathematical and source review.
\documentclass[11pt]{article}
\usepackage[margin=25mm]{geometry}
\usepackage{fontspec}
\setmainfont{Latin Modern Roman}
\usepackage{amsmath,amssymb,mathtools,unicode-math}
\setmathfont{Latin Modern Math}
\usepackage{enumitem,needspace,xurl,hyperref}
\setlength{\parindent}{0pt}
\setlength{\parskip}{5pt}
\setlength{\emergencystretch}{3em}
\newcommand{\N}{\mathbb{N}}
\newcommand{\Z}{\mathbb{Z}}
\newcommand{\Q}{\mathbb{Q}}
\newcommand{\R}{\mathbb{R}}
\newcommand{\C}{\mathbb{C}}
\newcommand{\F}{\mathbb{F}}
\newcommand{\A}{\mathbb{A}}
\newcommand{\PP}{\mathbb{P}}
\newcommand{\E}{\mathbb{E}}
\newcommand{\eps}{\varepsilon}
\newcommand{\dd}{\,\mathrm{d}}
\begin{document}
\section*{UnsolvedMath ID: 30004675}
\textbf{Source problem code:} catalogue entry
\subsection{Definitions and mathematical statement}
Let \(2^\omega=\{0,1\}^{\N}\) be Cantor space, where \(\N=\{0,1,2,\ldots\}\), and let \(A\subseteq 2^\omega\). The usual metric interpretation is the ultrametric \(d(x,x)=0\) and, for \(x\ne y\),
\[
 d(x,y)=2^{-\min\{n\in\N:x(n)\ne y(n)\}}.
\]
This metric convention is made explicit for interpretation; the supplied record does not specify a metric, and this convention requires source confirmation. Write \(\dim_H(A)\) and \(\dim_P(A)\) for the classical Hausdorff and packing dimensions of \(A\), respectively. In the terminology of this problem, \(A\) is \emph{regular} precisely when
\[
\dim_H(A)=\dim_P(A).
\]
A subset of Cantor space is \emph{analytic} if it is empty or is the image of Baire space \(\N^\N\) under a continuous map into \(2^\omega\).

The supplied record poses the following research questions.
\begin{enumerate}
\item For arbitrary \(A\subseteq 2^\omega\), can computability theory provide a partial characterization of the equality \(\dim_H(A)=\dim_P(A)\)? The record does not specify the computability property sought, allowable oracles, or the meaning of a partial characterization.
\item For analytic \(A\subseteq 2^\omega\) satisfying \(\dim_H(A)=\dim_P(A)\), can point-to-set principles provide an alternative proof of the referenced closed-subset result?
\item Can the analyticity hypothesis in that closed-subset result be removed?
\end{enumerate}
The closed-subset result is not stated in the supplied record. In particular, the required properties of the closed subset, the dimension relation it must satisfy, its quantifiers, and any additional hypotheses are missing. Thus the second and third questions cannot be converted into a complete quantified mathematical assertion from the supplied evidence. No particular closed-subset conclusion is asserted here.
\subsection{Short English statement}
The problem asks whether computability theory can help characterize subsets of Cantor space whose Hausdorff and packing dimensions agree, and whether point-to-set principles can reprove and extend a closed-subset theorem for regular analytic sets. The theorem itself and the intended computability criterion are unspecified, so those parts remain unresolved formulation issues.
\subsection{Sources}
\begin{itemize}
\item Local UnsolvedMath registry, permanent ID \(30004675\), title: \emph{Computable Characterizations of Hausdorff Packing Regularity}; exact source problem code: \texttt{catalogue entry}. The supplied record belongs to the Oberwolfach Reports Open Problems collection, gives proposed year \(2021\), and names no individual proposer.
\item Supplied source citation: \emph{Computability Theory (2022)}, Oberwolfach Reports, pp.\ 1149--1190, DOI \texttt{10.4171/owr/2021/21}. \url{https://doi.org/10.4171/owr/2021/21}. The background describes these reports as documenting mathematical workshops and their open-problem discussions. The citation's \(2022\) designation and the \(2021\) report identification are retained as supplied.
\item Stull, \emph{Optimal oracles for point-to-set principles}. Supplied literature source: \url{https://arxiv.org/abs/2101.11152}.
\item Peng, Wu and Yu, \emph{Some consequences of TD and sTD}. Supplied literature source: \url{https://doi.org/10.1017/jsl.2023.29}.
\item \emph{Oberwolfach Report 21/2021}. Supplied literature source: \url{https://ems.press/content/serial-article-files/46899}.
\item Literature assessment dated exactly \texttt{2026-08-22}: the supplied curation labels the problem partially solved and reports that point-to-set principles and optimal-oracle results settle substantial analytic and regular-set cases, while determinacy-based work extends closed-subset conclusions beyond analytic sets under extra axioms. It further reports that the full computability characterization and unconditional removal of analyticity remain open. This is supplied dated curation derived from the OWR record, not an independently verified literature assessment; its time-sensitive status and cited claims require independent verification.
\item The supplied statement status is \texttt{corrected\_verified}. Its accompanying verification note says that only formula delimiters, notation, grammar, and typography were repaired, with mathematical content and scope preserved. This supplied label does not establish mathematical correctness or resolve the missing formulation details.
\end{itemize}
\end{document}
